\documentclass[11pt, a4paper]{article}

\usepackage{fontspec}
\usepackage{geometry}
\geometry{a4paper, margin=2cm}
\usepackage{titlesec}
\usepackage{titling}
\usepackage{enumitem}
\setlist{nolistsep}
\usepackage{hyperref}
\usepackage{kotex}
\usepackage{array}
\usepackage{setspace}
\usepackage{xcolor}
\usepackage{enumitem}
\usepackage{fontawesome}
\definecolor{pastelgreen}{RGB}{120, 220, 120}
\definecolor{darkergreen}{RGB}{0, 128, 0}
\definecolor{darkeryellow}{RGB}{254, 176, 36}
\newcommand{\textbr}[1]{\textbf{\textcolor{bonusSteelBlue}{#1}}}
\newcommand{\textcb}[1]{\textcolor{bonusSteelBlue}{(#1)}}
\definecolor{bonusTeal}{RGB}{0, 128, 128}
\definecolor{bonusCoral}{RGB}{255, 127, 80}
\definecolor{bonusPlum}{RGB}{142, 69, 133}
\definecolor{bonusOlive}{RGB}{107, 142, 35}
\definecolor{bonusSteelBlue}{RGB}{70, 130, 180}
\definecolor{mediumGray}{RGB}{246,246,246}
\hypersetup{
    colorlinks=true,
    linkcolor=blue,
    filecolor=magenta,
    urlcolor=blue,
    pdftitle={Overleaf Example},
    pdfpagemode=FullScreen,
    }


\titleformat{\section}
{\large\bfseries}
{}
{0em}
{}[\titlerule]

\renewcommand{\maketitle}{
\begin{center}
{\huge\bfseries
\theauthor}

\vspace{.25em}
010-2805-0547 | tyeirdoo06@gmail.com
\end{center}
}

\begin{document}
\pagecolor{mediumGray}
\onehalfspacing


\title{이력서}
\author{두혁진}

\maketitle

\section*{경력}
\textbf{(재)경산이노베이션아카데미} (2024년 2월 - 현재)
\begin{itemize}
    \item \textbf{Prometheus \& Grafana 기반 모니터링 시스템 구축:} 워크스테이션 및 서버의 핵심 지표(CPU, Memory, Disk Usage 등)를 수집하고 Grafana 대시보드를 구축하여 장애 발생 시 평균 감지 시간(MTTD) 단축 및 장애 전파 사전 방지.
    \item \textbf{Rust 기반 플랫폼 운영 자동화 도구 개발:} 반복적인 장애 대응(세션 종료, 홈 리셋 등)을 자동화하는 API 서버(GsRoot)를 개발하여 24시간 안정적인 서비스 운영 기반을 마련하고 관리자의 수동 개입 감소에 기여.
    \item \textbf{데이터 수집 서버 개발 및 운영:} 교육생 데이터를 API와 크롤링으로 수집하고 Sqlite3에 저장하는 서버를 개발.
\end{itemize}\\

\textbf{42 서울 운영팀 (인턴)} (2023년 10월 - 2023년 12월)

업무: Ansible을 이용한 컴퓨터 클러스터 구성 관리, 42서울 내부 과제 명세 작성.

\section*{보유 기술 \textcolor{bonusSteelBlue}{(관심분야)}}
\begin{tabular}{r|p{0.7\textwidth}}
    \textbf{Language}& \textbf{Rust}, Bash\\
    \textbf{DB}& SQLite, PostgreSQL \\
    \textbf{VCS}& Git \\
    \textbf{Tooling \cdot~DevOps}& Ansible, GitHub \textbf{Actions}, Prometheus \& Grafana \textcb{(ELK)} \\
    \textbf{Editor}& Neovim
\end{tabular}

\section*{프로젝트}
\textbf{GsRoot (Rust)}

\textbf{배경}: 교육용 워크스테이션 사용 시 장애 해결을 위해 관리자 개입이 필요하며, 24시간 운영 환경에서 즉각적인 대응이 어려운 문제 발생.\\

\textbf{목표}: 교육용 워크스테이션에 대한 \textbf{서비스 안정성 확보} 및 \textbf{운영 자동화}를 위한 API 제공(재시작, 세션 종료, 홈 리셋).

\section*{개선 / 문제해결 사례}
\textbf{Humans of 42} (2023년 10월 - 11월)

개요: 유지보수가 중단된 기존 사이트를 Next.js 기반의 정적 페이지로 재구성.\\

\begin{tabular}{l p{0.8\textwidth}}
    1. 문제:& 사이트 인증서 만료로 인한 접속 불가 문제 반복\\
    2. 해결:& Next.js로 기술 스택을 현대화하고, \textbf{GitHub Actions 기반 CI/CD 파이프라인을 구축}하여 배포 자동화\\
    3. 결과:& 서버 비용 절감(월 5만원 \rightarrow 0원), 배포 자동화를 통한 운영 부담 해소, Github page의 인증서 자동 갱신으로 문제 해결
\end{tabular}

\textbf{} (2023년 10월 - 11월)

개요: 유지보수가 중단된 기존 사이트를 Next.js 기반의 정적 페이지로 재구성.\\

\begin{tabular}{l p{0.8\textwidth}}
    1. 문제:& 사이트 인증서 만료로 인한 접속 불가 문제 반복\\
    2. 해결:& Next.js로 기술 스택을 현대화하고, \textbf{GitHub Actions 기반 CI/CD 파이프라인을 구축}하여 배포 자동화\\
    3. 결과:& 서버 비용 절감(월 5만원 \rightarrow   0원), 배포 자동화를 통한 운영 부담 해소, Github page의 인증서 자동 갱신으로 문제 해결
\end{tabular}

\section*{도전 / 실패 사례}
\textbf{42 Region} (2022년 9월 - 12월)\\
개요: OpenStack을 활용한 프라이빗 클라우드 서비스 구축 시도.\\
\begin{tabular}{l p{0.8\textwidth}}
    1. 문제:& 인스턴스 생성 후 외부 인터넷 연결 불가 현상 발생\\
    2. 원인:& L1\textasciitilde L2 네트워크 계층에 대한 이해 부족으로 인한 Tagged VLAN 설정 미비\\
    3. 교훈:& 실패를 통해 인프라 지식의 중요성을 깨닫고 네트워크 기초 스터디를 주도. 이 경험을 계기로 컨테이너 오케스트레이션 기술(Kubernetes)과 퍼블릭 클라우드(AWS, GCP)에 대한 학습을 시작.
\end{tabular}

\end{document}
